\documentclass[11pt,a4paper]{article}
\usepackage[utf8]{inputenc}
\usepackage[T1]{fontenc}
\usepackage[catalan]{babel}
\usepackage{lmodern}
\usepackage[sfdefault,lining]{FiraSans}
\usepackage[a4paper,top=2.3cm,bottom=2.2cm,left=1.8cm,right=1.8cm,headheight=1.3cm,headsep=0.4cm,footskip=1cm]{geometry}
\usepackage{amsmath,amssymb}
\usepackage{icomma}
\usepackage{graphicx}
\usepackage[table]{xcolor}
\usepackage{tikz}
\usepackage{fontawesome5}
\usepackage{ccicons}
\usepackage{fancyhdr}
\usepackage{lastpage}
\usepackage{tasks}
\usepackage{needspace}
\usepackage[most]{tcolorbox}
\usepackage{array}
\usepackage[hidelinks,pdftitle={Fitxa d'exercicis de la unitat 11: Probabilitat. 4t d'ESO Matemàtiques A},pdfsubject={Matemàtiques A 4t ESO},pdfkeywords={probabilitat, experiment aleatori, espai mostral, esdeveniments, unió, intersecció, esdeveniment contrari, freqüència relativa, llei de Laplace, experiències compostes, experiències independents, experiències dependents, probabilitat condicionada, diagrama d'arbre, taula de contingència}]{hyperref}

\definecolor{navy}{HTML}{1B2748}
\definecolor{accent}{HTML}{E8603A}
\definecolor{teal}{HTML}{1F9E90}
\definecolor{soft}{HTML}{EEF1F8}
\definecolor{muted}{HTML}{5D6685}

\setlength{\parindent}{0pt}
\setlength{\parskip}{3pt}
\settasks{label=\alph*), label-format=\bfseries\color{navy}, label-width=1.3em, item-indent=1.8em, column-sep=0.6em, after-item-skip=5pt, before-skip=3pt, after-skip=4pt}

% ---------- capçalera i peu ----------
\newcommand{\llicencia}{\href{https://creativecommons.org/licenses/by-nc-sa/4.0/deed.ca}{\ccbyncsa}\ \ Llicència Creative Commons Reconeixement-NoComercial-CompartirIgual 4.0 Internacional (CC BY-NC-SA 4.0)}
\pagestyle{fancy}
\fancyhf{}
\fancyhead[L]{\raisebox{-0.35\height}{\includegraphics[height=0.95cm]{logo-gva.png}}}
\fancyhead[C]{\small\color{muted}Fitxa d'exercicis · Unitat 11: Probabilitat · 4t ESO A}
\fancyhead[R]{\raisebox{-0.35\height}{\includegraphics[height=0.95cm]{logo-vermellar.png}}}
\fancyfoot[L]{\scriptsize\color{muted}\llicencia}
\fancyfoot[R]{\small\color{muted}\thepage}
\renewcommand{\headrulewidth}{0.4pt}
\renewcommand{\headrule}{\color{muted!40}\hrule height \headrulewidth}
\fancypagestyle{primera}{\fancyhf{}\renewcommand{\headrulewidth}{0pt}
  \fancyfoot[L]{\scriptsize\color{muted}\llicencia}
  \fancyfoot[R]{\small\color{muted}\thepage}}

% ---------- nivells i exercicis ----------
\newcommand{\estrelles}[1]{{\color{accent}\ifcase#1\or\faStar\faStar[regular]\faStar[regular]\or\faStar\faStar\faStar[regular]\or\faStar\faStar\faStar\fi}}
\newcounter{exer}
\newcommand{\ex}[1]{\par\needspace{4\baselineskip}\medskip\stepcounter{exer}\noindent{\color{navy}\bfseries\theexer.}\ \ {\small\estrelles{#1}}\ \ }

\newcommand{\seccio}[2]{\par\needspace{10\baselineskip}\bigskip
  \begin{tcolorbox}[colback=navy,colframe=navy,coltext=white,arc=2mm,boxrule=0pt,left=3mm,right=3mm,top=1.2mm,bottom=1.2mm]
  {\large\bfseries #1\quad #2}
  \end{tcolorbox}\vspace{-2pt}}

\newcolumntype{C}[1]{>{\centering\arraybackslash}p{#1}}
\usepackage{pifont}
\usetikzlibrary{trees}
% quadre de recordatori (només a la fitxa d'exercicis)
\newcommand{\recorda}[1]{\par\smallskip\begin{tcolorbox}[colback=accent!7,colframe=accent!55,arc=2mm,boxrule=0.5pt,left=2.5mm,right=2.5mm,top=1mm,bottom=1mm]\small{\color{accent}\faLightbulb[regular]}\ \textbf{Recorda:}\ #1\end{tcolorbox}}
\newcommand{\res}[1]{#1}\newcommand{\rest}[1]{#1}

\def\fmneed{14}
% =====================================================================
% Capçalera FentMates de les fitxes: un color per a cada curs.
% Els colors són els de la pàgina inicial (LEVELS a assets/js/materials.js);
% el logo ve de _eines/logo/pdf/ (python3 _eines/logo/genera.py).
%
% Ús, al final del preàmbul de la unitat:
%   % =====================================================================
% Capçalera FentMates de les fitxes: un color per a cada curs.
% Els colors són els de la pàgina inicial (LEVELS a assets/js/materials.js);
% el logo ve de _eines/logo/pdf/ (python3 _eines/logo/genera.py).
%
% Ús, al final del preàmbul de la unitat:
%   % =====================================================================
% Capçalera FentMates de les fitxes: un color per a cada curs.
% Els colors són els de la pàgina inicial (LEVELS a assets/js/materials.js);
% el logo ve de _eines/logo/pdf/ (python3 _eines/logo/genera.py).
%
% Ús, al final del preàmbul de la unitat:
%   \input{../comu/fentmates.tex}
%   \fmcurs{1eso}
% i, al començament del document:
%   \capcalera{Matemàtiques}{1r ESO}{Fitxa d'exercicis · Unitat 2}{Divisibilitat}
% Opcional abans de \input: \def\fmneed{14} (espai mínim abans de cada secció).
% =====================================================================
\makeatletter
\def\fm@logodir{../../logo/pdf}
\providecommand{\fmneed}{10}
\def\fm@cursnom{}

% clau del curs, color i color del text damunt del color
\newcommand{\fm@curs}[3]{\@namedef{fm@color@#1}{#2}\@namedef{fm@text@#1}{#3}}
\fm@curs{1eso}{E6B84A}{navy}
\fm@curs{2eso}{3B82F6}{white}
\fm@curs{3eso}{FF7A50}{white}
\fm@curs{4eso-a}{A78BFA}{white}
\fm@curs{4eso-b}{6D28D9}{white}
\fm@curs{1batx-mi}{9B1C1C}{white}
\fm@curs{1batx-ccss}{EF6B6B}{white}
\fm@curs{2batx-mii}{1E6B47}{white}
\fm@curs{2batx-ccss}{2BB3A3}{white}

\newcommand{\fmcurs}[1]{%
  \@ifundefined{fm@color@#1}{\PackageError{fentmates}{Curs desconegut: #1}{}}{}%
  \edef\fm@tmp{\noexpand\definecolor{curs}{HTML}{\@nameuse{fm@color@#1}}%
               \noexpand\colorlet{curstext}{\@nameuse{fm@text@#1}}}\fm@tmp
  \colorlet{cursfosc}{curs!72!black}%
  \def\fm@logo{\fm@logodir/vertical-#1.pdf}%
  \def\fm@marca{\fm@logodir/marca-#1.pdf}}

% ---------- capçalera de la primera pàgina ----------
% #1 matèria (petit), #2 curs, #3 tipus de document i unitat, #4 títol
\newcommand{\capcalera}[4]{%
  \gdef\fm@cursnom{#2}%
  \par\noindent
  \begin{tcolorbox}[colback=curs,colframe=curs,boxrule=0pt,arc=2.2mm,width=3.5cm,height=2.15cm,
      valign=center,left=3mm,right=2mm,top=1mm,bottom=1mm,nobeforeafter,box align=center]
    \color{curstext}{\footnotesize #1}\par\vspace{1pt}{\fontsize{19}{21}\selectfont\bfseries #2\par}%
  \end{tcolorbox}\hspace{5mm}%
  \begin{minipage}[c]{\dimexpr\linewidth-3.5cm-5mm-2.9cm\relax}
    {\color{cursfosc}\small\bfseries\MakeUppercase{#3}}\par\vspace{2pt}
    {\color{navy}\fontsize{20}{23}\selectfont\bfseries #4\par}
  \end{minipage}\hfill
  \begin{minipage}[c]{2.6cm}\raggedleft\includegraphics[width=2.6cm]{\fm@logo}\end{minipage}\par
  \vspace{2pt}}

% ---------- barres de secció del color del curs ----------
\renewcommand{\seccio}[2]{\par\needspace{\fmneed\baselineskip}\bigskip
  \begin{tcolorbox}[colback=curs,colframe=curs,coltext=curstext,arc=2mm,boxrule=0pt,left=3mm,right=3mm,top=1.2mm,bottom=1.2mm]
  {\large\bfseries #1\quad #2}
  \end{tcolorbox}\vspace{-2pt}}

% ---------- estrelles de dificultat del color del curs (a Batxillerat no se'n posen) ----------
\@ifundefined{estrelles}{}{%
  \renewcommand{\estrelles}[1]{{\color{curs}\ifcase#1\or\faStar\faStar[regular]\faStar[regular]\or\faStar\faStar\faStar[regular]\or\faStar\faStar\faStar\fi}}}

% ---------- capçalera i peu de les pàgines ----------
% primera pàgina: logos institucionals al peu
\fancypagestyle{primera}{\fancyhf{}\renewcommand{\headrulewidth}{0pt}
  \fancyfoot[L]{\raisebox{-0.32\height}{\includegraphics[height=0.78cm]{logo-gva.png}}\hspace{5mm}%
    \raisebox{-0.32\height}{\includegraphics[height=0.78cm]{logo-vermellar.png}}}
  \fancyfoot[R]{\scriptsize\color{muted}\href{https://creativecommons.org/licenses/by-nc-sa/4.0/deed.ca}{\ccbyncsa}\ \ CC BY-NC-SA 4.0 · fentmates.org\qquad\small\thepage}}
% la resta: el símbol a l'esquerra i el curs a la dreta, amb el títol al centre com abans
\fancyhead[L]{\raisebox{-0.3\height}{\includegraphics[height=0.62cm]{\fm@marca}}}
\fancyhead[R]{\small\bfseries\color{cursfosc}\fm@cursnom}
\renewcommand{\headrule}{\color{curs!45}\hrule height \headrulewidth}
\makeatother

%   \fmcurs{1eso}
% i, al començament del document:
%   \capcalera{Matemàtiques}{1r ESO}{Fitxa d'exercicis · Unitat 2}{Divisibilitat}
% Opcional abans de \input: \def\fmneed{14} (espai mínim abans de cada secció).
% =====================================================================
\makeatletter
\def\fm@logodir{../../logo/pdf}
\providecommand{\fmneed}{10}
\def\fm@cursnom{}

% clau del curs, color i color del text damunt del color
\newcommand{\fm@curs}[3]{\@namedef{fm@color@#1}{#2}\@namedef{fm@text@#1}{#3}}
\fm@curs{1eso}{E6B84A}{navy}
\fm@curs{2eso}{3B82F6}{white}
\fm@curs{3eso}{FF7A50}{white}
\fm@curs{4eso-a}{A78BFA}{white}
\fm@curs{4eso-b}{6D28D9}{white}
\fm@curs{1batx-mi}{9B1C1C}{white}
\fm@curs{1batx-ccss}{EF6B6B}{white}
\fm@curs{2batx-mii}{1E6B47}{white}
\fm@curs{2batx-ccss}{2BB3A3}{white}

\newcommand{\fmcurs}[1]{%
  \@ifundefined{fm@color@#1}{\PackageError{fentmates}{Curs desconegut: #1}{}}{}%
  \edef\fm@tmp{\noexpand\definecolor{curs}{HTML}{\@nameuse{fm@color@#1}}%
               \noexpand\colorlet{curstext}{\@nameuse{fm@text@#1}}}\fm@tmp
  \colorlet{cursfosc}{curs!72!black}%
  \def\fm@logo{\fm@logodir/vertical-#1.pdf}%
  \def\fm@marca{\fm@logodir/marca-#1.pdf}}

% ---------- capçalera de la primera pàgina ----------
% #1 matèria (petit), #2 curs, #3 tipus de document i unitat, #4 títol
\newcommand{\capcalera}[4]{%
  \gdef\fm@cursnom{#2}%
  \par\noindent
  \begin{tcolorbox}[colback=curs,colframe=curs,boxrule=0pt,arc=2.2mm,width=3.5cm,height=2.15cm,
      valign=center,left=3mm,right=2mm,top=1mm,bottom=1mm,nobeforeafter,box align=center]
    \color{curstext}{\footnotesize #1}\par\vspace{1pt}{\fontsize{19}{21}\selectfont\bfseries #2\par}%
  \end{tcolorbox}\hspace{5mm}%
  \begin{minipage}[c]{\dimexpr\linewidth-3.5cm-5mm-2.9cm\relax}
    {\color{cursfosc}\small\bfseries\MakeUppercase{#3}}\par\vspace{2pt}
    {\color{navy}\fontsize{20}{23}\selectfont\bfseries #4\par}
  \end{minipage}\hfill
  \begin{minipage}[c]{2.6cm}\raggedleft\includegraphics[width=2.6cm]{\fm@logo}\end{minipage}\par
  \vspace{2pt}}

% ---------- barres de secció del color del curs ----------
\renewcommand{\seccio}[2]{\par\needspace{\fmneed\baselineskip}\bigskip
  \begin{tcolorbox}[colback=curs,colframe=curs,coltext=curstext,arc=2mm,boxrule=0pt,left=3mm,right=3mm,top=1.2mm,bottom=1.2mm]
  {\large\bfseries #1\quad #2}
  \end{tcolorbox}\vspace{-2pt}}

% ---------- estrelles de dificultat del color del curs (a Batxillerat no se'n posen) ----------
\@ifundefined{estrelles}{}{%
  \renewcommand{\estrelles}[1]{{\color{curs}\ifcase#1\or\faStar\faStar[regular]\faStar[regular]\or\faStar\faStar\faStar[regular]\or\faStar\faStar\faStar\fi}}}

% ---------- capçalera i peu de les pàgines ----------
% primera pàgina: logos institucionals al peu
\fancypagestyle{primera}{\fancyhf{}\renewcommand{\headrulewidth}{0pt}
  \fancyfoot[L]{\raisebox{-0.32\height}{\includegraphics[height=0.78cm]{logo-gva.png}}\hspace{5mm}%
    \raisebox{-0.32\height}{\includegraphics[height=0.78cm]{logo-vermellar.png}}}
  \fancyfoot[R]{\scriptsize\color{muted}\href{https://creativecommons.org/licenses/by-nc-sa/4.0/deed.ca}{\ccbyncsa}\ \ CC BY-NC-SA 4.0 · fentmates.org\qquad\small\thepage}}
% la resta: el símbol a l'esquerra i el curs a la dreta, amb el títol al centre com abans
\fancyhead[L]{\raisebox{-0.3\height}{\includegraphics[height=0.62cm]{\fm@marca}}}
\fancyhead[R]{\small\bfseries\color{cursfosc}\fm@cursnom}
\renewcommand{\headrule}{\color{curs!45}\hrule height \headrulewidth}
\makeatother

%   \fmcurs{1eso}
% i, al començament del document:
%   \capcalera{Matemàtiques}{1r ESO}{Fitxa d'exercicis · Unitat 2}{Divisibilitat}
% Opcional abans de \input: \def\fmneed{14} (espai mínim abans de cada secció).
% =====================================================================
\makeatletter
\def\fm@logodir{../../logo/pdf}
\providecommand{\fmneed}{10}
\def\fm@cursnom{}

% clau del curs, color i color del text damunt del color
\newcommand{\fm@curs}[3]{\@namedef{fm@color@#1}{#2}\@namedef{fm@text@#1}{#3}}
\fm@curs{1eso}{E6B84A}{navy}
\fm@curs{2eso}{3B82F6}{white}
\fm@curs{3eso}{FF7A50}{white}
\fm@curs{4eso-a}{A78BFA}{white}
\fm@curs{4eso-b}{6D28D9}{white}
\fm@curs{1batx-mi}{9B1C1C}{white}
\fm@curs{1batx-ccss}{EF6B6B}{white}
\fm@curs{2batx-mii}{1E6B47}{white}
\fm@curs{2batx-ccss}{2BB3A3}{white}

\newcommand{\fmcurs}[1]{%
  \@ifundefined{fm@color@#1}{\PackageError{fentmates}{Curs desconegut: #1}{}}{}%
  \edef\fm@tmp{\noexpand\definecolor{curs}{HTML}{\@nameuse{fm@color@#1}}%
               \noexpand\colorlet{curstext}{\@nameuse{fm@text@#1}}}\fm@tmp
  \colorlet{cursfosc}{curs!72!black}%
  \def\fm@logo{\fm@logodir/vertical-#1.pdf}%
  \def\fm@marca{\fm@logodir/marca-#1.pdf}}

% ---------- capçalera de la primera pàgina ----------
% #1 matèria (petit), #2 curs, #3 tipus de document i unitat, #4 títol
\newcommand{\capcalera}[4]{%
  \gdef\fm@cursnom{#2}%
  \par\noindent
  \begin{tcolorbox}[colback=curs,colframe=curs,boxrule=0pt,arc=2.2mm,width=3.5cm,height=2.15cm,
      valign=center,left=3mm,right=2mm,top=1mm,bottom=1mm,nobeforeafter,box align=center]
    \color{curstext}{\footnotesize #1}\par\vspace{1pt}{\fontsize{19}{21}\selectfont\bfseries #2\par}%
  \end{tcolorbox}\hspace{5mm}%
  \begin{minipage}[c]{\dimexpr\linewidth-3.5cm-5mm-2.9cm\relax}
    {\color{cursfosc}\small\bfseries\MakeUppercase{#3}}\par\vspace{2pt}
    {\color{navy}\fontsize{20}{23}\selectfont\bfseries #4\par}
  \end{minipage}\hfill
  \begin{minipage}[c]{2.6cm}\raggedleft\includegraphics[width=2.6cm]{\fm@logo}\end{minipage}\par
  \vspace{2pt}}

% ---------- barres de secció del color del curs ----------
\renewcommand{\seccio}[2]{\par\needspace{\fmneed\baselineskip}\bigskip
  \begin{tcolorbox}[colback=curs,colframe=curs,coltext=curstext,arc=2mm,boxrule=0pt,left=3mm,right=3mm,top=1.2mm,bottom=1.2mm]
  {\large\bfseries #1\quad #2}
  \end{tcolorbox}\vspace{-2pt}}

% ---------- estrelles de dificultat del color del curs (a Batxillerat no se'n posen) ----------
\@ifundefined{estrelles}{}{%
  \renewcommand{\estrelles}[1]{{\color{curs}\ifcase#1\or\faStar\faStar[regular]\faStar[regular]\or\faStar\faStar\faStar[regular]\or\faStar\faStar\faStar\fi}}}

% ---------- capçalera i peu de les pàgines ----------
% primera pàgina: logos institucionals al peu
\fancypagestyle{primera}{\fancyhf{}\renewcommand{\headrulewidth}{0pt}
  \fancyfoot[L]{\raisebox{-0.32\height}{\includegraphics[height=0.78cm]{logo-gva.png}}\hspace{5mm}%
    \raisebox{-0.32\height}{\includegraphics[height=0.78cm]{logo-vermellar.png}}}
  \fancyfoot[R]{\scriptsize\color{muted}\href{https://creativecommons.org/licenses/by-nc-sa/4.0/deed.ca}{\ccbyncsa}\ \ CC BY-NC-SA 4.0 · fentmates.org\qquad\small\thepage}}
% la resta: el símbol a l'esquerra i el curs a la dreta, amb el títol al centre com abans
\fancyhead[L]{\raisebox{-0.3\height}{\includegraphics[height=0.62cm]{\fm@marca}}}
\fancyhead[R]{\small\bfseries\color{cursfosc}\fm@cursnom}
\renewcommand{\headrule}{\color{curs!45}\hrule height \headrulewidth}
\makeatother

\fmcurs{4eso-a}

\begin{document}
\thispagestyle{primera}

\capcalera{Matemàtiques A}{4t ESO}{Fitxa d'exercicis · Unitat 11}{Probabilitat}
\vspace{10pt}
Nom i cognoms: \rule{8.4cm}{0.4pt}\hfill Grup: \rule{1.6cm}{0.4pt}\hfill Data: \rule{2.2cm}{0.4pt}
\vspace{8pt}

\begin{tcolorbox}[colback=soft,colframe=soft,arc=2mm,left=3mm,right=3mm,top=1.5mm,bottom=1.5mm]
\small Exercicis graduats per nivells de dificultat:\quad \estrelles{1}\ bàsic \qquad \estrelles{2}\ mitjà \qquad \estrelles{3}\ avançat.\\[2pt] Resol-los al quadern mostrant tots els passos: escriu l'espai mostral, fes un diagrama d'arbre o una taula quan calga i dona les probabilitats en forma de fracció irreductible (o en decimal, si les dades ho són). Les seccions segueixen l'ordre dels apunts de la unitat (11.1--11.7).
\end{tcolorbox}

% =====================================================
\seccio{1}{Experiments}

\ex{1} Indica si són experiments aleatoris o deterministes:\begin{tasks}(2)
  \task Llançar un dau i anotar el resultat.
  \task Mesurar la temperatura a què bull l'aigua al nivell de la mar.
  \task Traure una carta d'una baralla sense mirar.
  \task Calcular quant pagarem per 3 kg de pomes a 2 €/kg.
  \task Endevinar el número guanyador de la loteria.
  \task Deixar caure una pilota i mirar si cau a terra.
\end{tasks}\recorda{Un experiment és \textbf{aleatori} si no podem predir-ne el resultat, encara que el repetim en les mateixes condicions. És \textbf{determinista} si el resultat és sempre el mateix i el podem saber abans de fer-lo.}

\ex{1} En una bossa hi ha 9 boles roges i 1 de blava. Traiem una bola sense mirar. Descriu cada esdeveniment amb una d'aquestes paraules: \emph{segur}, \emph{molt probable}, \emph{poc probable}, \emph{impossible}.\par\smallskip\begin{center}\begin{tikzpicture}[x=0.42cm,y=0.42cm,baseline=(current bounding box.center)]
\draw[muted,thick,fill=soft] (-0.2,0.1) -- (-0.2,-1.70) to[out=-90,in=180] (2.50,-2.30) to[out=0,in=-90] (5.20,-1.70) -- (5.20,0.1);
\draw[muted,line width=1.6pt,line cap=round] (-0.4,0.1) -- (5.40,0.1);
\filldraw[draw=black!50,fill=red!75!black] (0.50,-0.45) circle (0.38);
\filldraw[draw=black!50,fill=red!75!black] (1.50,-0.45) circle (0.38);
\filldraw[draw=black!50,fill=red!75!black] (2.50,-0.45) circle (0.38);
\filldraw[draw=black!50,fill=red!75!black] (3.50,-0.45) circle (0.38);
\filldraw[draw=black!50,fill=red!75!black] (4.50,-0.45) circle (0.38);
\filldraw[draw=black!50,fill=red!75!black] (0.50,-1.45) circle (0.38);
\filldraw[draw=black!50,fill=red!75!black] (1.50,-1.45) circle (0.38);
\filldraw[draw=black!50,fill=red!75!black] (2.50,-1.45) circle (0.38);
\filldraw[draw=black!50,fill=red!75!black] (3.50,-1.45) circle (0.38);
\filldraw[draw=black!50,fill=blue!70!black] (4.50,-1.45) circle (0.38);
\end{tikzpicture}\end{center}\begin{tasks}(4)
  \task Traure roja.
  \task Traure blava.
  \task Traure verda.
  \task Traure roja o blava.
\end{tasks}

\ex{1} Escriu l'espai mostral d'aquests experiments:\begin{tasks}(2)
  \task Llançar una moneda.
  \task Llançar un dau de quatre cares numerades de l'1 al 4.
  \task Girar una ruleta de 5 sectors iguals (roig, blau, verd, groc i negre).
  \task Traure una bola d'una bossa amb boles roges i blaves.
  \task Triar a l'atzar un dia de la setmana.
  \task Preguntar a algú quantes germanes té (com a màxim, 3).
\end{tasks}\recorda{L'\textbf{espai mostral} $E$ és el conjunt de tots els resultats possibles, i un \textbf{esdeveniment} és un subconjunt d'$E$ (un grup de resultats).}

\ex{2} En Marc diu que llançar un dau és un experiment determinista, «perquè ja sabem que eixirà un nombre de l'1 al 6». Té raó? Explica-ho.

\ex{2} Llancem dues monedes. Fes un diagrama d'arbre i escriu l'espai mostral i els esdeveniments $A=$ «almenys una cara», $B=$ «cap cara» i $C=$ «les dues iguals».

\ex{2} Una família té tres fills. Anotem si cadascun és xiqueta (A) o xiquet (O), per ordre de naixement. Fes un diagrama d'arbre, escriu l'espai mostral i l'esdeveniment «exactament dues xiquetes».

\ex{2} Un menú té dos primers plats, sopa (S) o amanida (A), i tres segons, carn (C), peix (P) o truita (T). Triem el menú a l'atzar. Fes un diagrama d'arbre, escriu l'espai mostral i l'esdeveniment «el menú porta peix».

\ex{3} Formem nombres de tres xifres diferents amb les xifres 1, 2 i 3. Escriu l'espai mostral i els esdeveniments «el nombre és parell» i «el nombre és més gran que 200».


% =====================================================
\seccio{2}{Esdeveniments}

\ex{1} Traiem una bola d'una urna amb 12 boles numerades de l'1 al 12. Escriu aquests esdeveniments i digues de quin tipus és cadascun (elemental, compost, segur o impossible):\begin{tasks}(3)
  \task $A=$ «parell»
  \task $B=$ «múltiple de 5»
  \task $C=$ «més gran que 12»
  \task $D=$ «divisor de 12»
  \task $F=$ «menor que 13»
  \task $G=$ «múltiple d'11»
\end{tasks}

\ex{1} Llancem un dau de 8 cares numerades de l'1 al 8. Escriu l'espai mostral i els esdeveniments:\begin{tasks}(4)
  \task «parell»
  \task «primer»
  \task «més gran que 5»
  \task «divisor de 8»
\end{tasks}

\ex{1} Llancem un dau. Siguen $A=$ «múltiple de 3» i $B=$ «menor que 3». Calcula $A\cup B$, $A\cap B$, $\overline{A}$ i $\overline{B}$. Són incompatibles? Són contraris?\recorda{$A\cup B$ («$A$ o $B$»): elements d'$A$ o de $B$. \ $A\cap B$ («$A$ i $B$»): elements comuns. \ $\overline{A}$ («no $A$»): elements de $E$ que no són en $A$. Si $A\cap B=\varnothing$, són \textbf{incompatibles}.}

\ex{1} En una urna hi ha deu boles numerades de l'1 al 10. Siguen $A=$ «parell» i $B=$ «primer». Calcula $A\cup B$, $A\cap B$, $\overline{A}$ i $\overline{B}$.

\ex{2} Traiem una carta d'una baralla espanyola de 40 cartes. Siguen $O=$ «oros» i $A=$ «as». Quantes cartes té cada un d'aquests esdeveniments: $O\cap A$, $O\cup A$ i $\overline{O}$?

\ex{2} Llancem tres monedes. Siguen $A=$ «almenys dues cares» i $B=$ «la primera és creu». Escriu $A$, $B$, $A\cap B$, $A\cup B$ i $\overline{A}$.

\ex{3} Traiem una bola d'una urna amb boles de l'1 al 12. Siguen $A=$ «múltiple de 2» i $B=$ «múltiple de 3». Escriu amb operacions i calcula:
\begin{tasks}(1)
  \task «és múltiple de 2 però no de 3»
  \task «no és múltiple ni de 2 ni de 3». Comprova que b) és el contrari d'$A\cup B$
\end{tasks}


% =====================================================
\seccio{3}{Probabilitat d'un esdeveniment}

\ex{1} Hem llançat un dau 50 vegades i hem obtingut aquests resultats. Completa la taula amb les freqüències relatives (en decimal) i calcula la freqüència relativa de «traure parell».\par\smallskip\begin{center}{\small\renewcommand{\arraystretch}{1.2}\begin{tabular}{l|cccccc}
\rowcolor{soft} Resultat & 1 & 2 & 3 & 4 & 5 & 6\\\hline
Freq. absoluta & 7 & 9 & 8 & 10 & 6 & 10\\
Freq. relativa &  &  &  &  &  & \\
\end{tabular}}\end{center}\recorda{Si un esdeveniment ocorre $f$ vegades en $n$ repeticions, la seua \textbf{freqüència relativa} és $h=\frac{f}{n}$.}

\ex{1} Quins d'aquests nombres poden ser la probabilitat d'un esdeveniment?\begin{tasks}(4)
  \task $0{,}7$
  \task $1{,}2$
  \task $-0{,}1$
  \task $\frac32$
  \task $0$
  \task $100\,\%$
  \task $\frac54$
  \task $\frac{99}{100}$
\end{tasks}

\ex{1} Calcula la probabilitat de l'esdeveniment contrari:\begin{tasks}(3)
  \task $P(A)=0{,}62$
  \task $P(B)=\frac37$
  \task $P(C)=15\,\%$
\end{tasks}\recorda{$P(\overline{A})=1-P(A)$ \qquad $P(A\cup B)=P(A)+P(B)-P(A\cap B)$ \qquad Si són incompatibles, $P(A\cup B)=P(A)+P(B)$.}

\ex{2} Hem llançat una xinxeta diverses vegades i hem anotat quantes ha caigut amb la punta cap amunt. Completa la taula i estima la probabilitat que caiga amb la punta cap amunt.\par\smallskip\begin{center}{\small\renewcommand{\arraystretch}{1.2}\begin{tabular}{l|cccc}
\rowcolor{soft} Llançaments & 20 & 100 & 500 & 1000\\\hline
Punta amunt & 13 & 58 & 297 & 601\\
Freq. relativa &  &  &  & \\
\end{tabular}}\end{center}\recorda{\textbf{Llei dels grans nombres}: quan repetim un experiment moltes vegades, la freqüència relativa d'un esdeveniment s'estabilitza al voltant de la seua probabilitat.}

\ex{2} Un jugador d'handbol ha marcat 45 dels 60 penals que ha llançat. Estima la probabilitat que marque un penal. Quants penals esperes que marque dels 20 pròxims?

\ex{2} Un dau està trucat. Les probabilitats de l'1 al 5 són 0,15, 0,1, 0,2, 0,15 i 0,1. Calcula $P(6)$ i la probabilitat de traure parell.

\ex{2} Siguen $A$ i $B$ dos esdeveniments amb $P(A)=0{,}4$, $P(B)=0{,}5$ i $P(A\cap B)=0{,}15$. Calcula $P(A\cup B)$, la probabilitat que no ocórrega cap dels dos i la probabilitat que ocórrega $A$ però no $B$.

\ex{2} 
\begin{tasks}(1)
  \task $A$ i $B$ són incompatibles, amb $P(A)=0{,}3$ i $P(B)=0{,}45$. Calcula $P(A\cup B)$.
  \task Poden ser incompatibles dos esdeveniments amb $P(A)=0{,}6$ i $P(B)=0{,}5$? Per què?
\end{tasks}

\ex{3} En una bossa hi ha 30 boles, roges i blanques. Traiem una bola, n'anotem el color i la tornem a la bossa. Després de 300 extraccions n'han eixit 120 de roges. Quantes boles roges deu haver-hi?

\ex{3} Sabem que $P(A\cup B)=0{,}9$, $P(\overline{A})=0{,}3$ i $P(\overline{B})=0{,}4$. Calcula $P(A)$, $P(B)$ i $P(A\cap B)$. Són incompatibles?


% =====================================================
\seccio{4}{Càlcul de probabilitats}

\ex{1} Llancem un dau. Calcula la probabilitat de:\begin{tasks}(3)
  \task traure un 5
  \task traure parell
  \task traure més de 2
  \task traure un múltiple de 3
  \task traure un 7
  \task traure menys de 7
\end{tasks}\recorda{En un experiment \textbf{regular} (resultats equiprobables): \ $P(A)=\dfrac{\text{nombre de casos favorables}}{\text{nombre de casos possibles}}$.}

\ex{1} En una bossa hi ha 6 boles roges, 4 de blaves i 5 de verdes. Traiem una bola a l'atzar. Calcula $P(\text{roja})$, $P(\text{blava})$, $P(\text{verda})$ i $P(\text{no verda})$.\par\smallskip\begin{center}\begin{tikzpicture}[x=0.42cm,y=0.42cm,baseline=(current bounding box.center)]
\draw[muted,thick,fill=soft] (-0.2,0.1) -- (-0.2,-1.70) to[out=-90,in=180] (4.00,-2.30) to[out=0,in=-90] (8.20,-1.70) -- (8.20,0.1);
\draw[muted,line width=1.6pt,line cap=round] (-0.4,0.1) -- (8.40,0.1);
\filldraw[draw=black!50,fill=red!75!black] (0.50,-0.45) circle (0.38);
\filldraw[draw=black!50,fill=red!75!black] (1.50,-0.45) circle (0.38);
\filldraw[draw=black!50,fill=red!75!black] (2.50,-0.45) circle (0.38);
\filldraw[draw=black!50,fill=red!75!black] (3.50,-0.45) circle (0.38);
\filldraw[draw=black!50,fill=red!75!black] (4.50,-0.45) circle (0.38);
\filldraw[draw=black!50,fill=red!75!black] (5.50,-0.45) circle (0.38);
\filldraw[draw=black!50,fill=blue!70!black] (6.50,-0.45) circle (0.38);
\filldraw[draw=black!50,fill=blue!70!black] (7.50,-0.45) circle (0.38);
\filldraw[draw=black!50,fill=blue!70!black] (0.50,-1.45) circle (0.38);
\filldraw[draw=black!50,fill=blue!70!black] (1.50,-1.45) circle (0.38);
\filldraw[draw=black!50,fill=teal] (2.50,-1.45) circle (0.38);
\filldraw[draw=black!50,fill=teal] (3.50,-1.45) circle (0.38);
\filldraw[draw=black!50,fill=teal] (4.50,-1.45) circle (0.38);
\filldraw[draw=black!50,fill=teal] (5.50,-1.45) circle (0.38);
\filldraw[draw=black!50,fill=teal] (6.50,-1.45) circle (0.38);
\end{tikzpicture}\end{center}

\ex{1} Llancem dos daus i sumem els punts. Completa la taula de doble entrada i calcula la probabilitat que la suma siga:
\begin{tasks}(1)
  \task 8
  \task 2
  \task 10 o més
  \task senar
\end{tasks}

\ex{2} Traiem una carta d'una baralla espanyola de 40 cartes (oros, copes, espases i bastos, amb l'1 al 7, la sota, el cavall i el rei de cada pal). Calcula la probabilitat que siga:\begin{tasks}(3)
  \task de copes
  \task un as
  \task una figura d'espases
  \task una carta que no siga figura
  \task un as o un oro
  \task el cavall de bastos
\end{tasks}

\ex{2} Escrivim cada lletra de la paraula MATEMATIQUES en una targeta i en traiem una a l'atzar. Quina és la probabilitat que siga una vocal? I una M? I una consonant?

\ex{2} En Pau i la Maria juguen amb un dau: en Pau guanya si ix 5 o 6, i la Maria si ix 1, 2, 3 o 4. És just el joc? Com el canviaries perquè ho fóra?

\ex{2} Llancem dos daus. Calcula la probabilitat que:
\begin{tasks}(1)
  \task traguen el mateix nombre
  \task la diferència entre el major i el menor siga 2
  \task el producte siga parell
\end{tasks}

\ex{3} En una bossa hi ha 5 boles roges i unes quantes de blaves. Si la probabilitat de traure una bola blava és $\frac23$, quantes boles blaves hi ha?

\ex{3} Aquesta ruleta té tres sectors de 180°, 120° i 60°. És un experiment regular? Calcula la probabilitat de cada color.\par\smallskip\begin{center}\begin{tikzpicture}[scale=0.75,baseline=(current bounding box.center)]\fill[red!60] (0,0) -- (0:1.3) arc[start angle=0,end angle=180,radius=1.3] -- cycle;\fill[blue!55] (0,0) -- (180:1.3) arc[start angle=180,end angle=300,radius=1.3] -- cycle;\fill[teal!60] (0,0) -- (300:1.3) arc[start angle=300,end angle=360,radius=1.3] -- cycle;\draw[thick] (0,0) circle (1.3);\draw (0,0)--(0:1.3) (0,0)--(180:1.3) (0,0)--(300:1.3);\node[font=\scriptsize\bfseries] at (90:0.7) {180°};\node[font=\scriptsize\bfseries,white] at (240:0.75) {120°};\node[font=\scriptsize\bfseries] at (330:0.8) {60°};\end{tikzpicture}\end{center}


% =====================================================
\seccio{5}{Experiències compostes. Probabilitat}

\ex{1} Llancem una moneda i un dau. Fes el diagrama d'arbre. Quants resultats hi ha? Calcula la probabilitat de:
\begin{tasks}(1)
  \task cara i un 3
  \task creu i un nombre parell
  \task un nombre primer (amb cara o creu)
\end{tasks}
\recorda{En un \textbf{diagrama d'arbre}, cada camí és un resultat. Si posem probabilitats a les branques, la probabilitat d'un camí és el \textbf{producte} de les seues branques; si un esdeveniment té diversos camins, les \textbf{sumem}.}

\ex{1} Llancem dues monedes. Fes el diagrama d'arbre amb les probabilitats i calcula la probabilitat de traure dues cares i la de traure almenys una creu.

\ex{1} La probabilitat que ploga dissabte és 0,3 i la que ploga diumenge, 0,6 (suposem que són independents). Fes el diagrama d'arbre i calcula la probabilitat que:
\begin{tasks}(1)
  \task ploga els dos dies
  \task no ploga cap dia
  \task ploga almenys un dia
  \task ploga exactament un dia
\end{tasks}
\recorda{Si dues experiències són \textbf{independents}, $P(A\cap B)=P(A)\cdot P(B)$. Per a «almenys un», sovint és més ràpid el contrari: $P(\text{almenys un})=1-P(\text{cap})$.}

\ex{2} Llancem tres monedes. Calcula la probabilitat de:
\begin{tasks}(1)
  \task tres cares
  \task exactament una cara
  \task almenys dues cares
\end{tasks}

\ex{2} Una jugadora de futbol marca el 80 \% dels penals. En llança dos. Fes el diagrama d'arbre i calcula la probabilitat que marque tots dos, que no en marque cap i que en marque exactament un.

\ex{2} En una bossa hi ha 3 boles blanques (B) i 2 de negres (N). En traiem dues \textbf{amb reemplaçament}. Fes el diagrama d'arbre i calcula la probabilitat que:
\begin{tasks}(1)
  \task les dues siguen blanques
  \task les dues siguen negres
  \task n'isca una de cada color
\end{tasks}
\recorda{\textbf{Amb reemplaçament} tornem la bola a la bossa: la segona extracció té les mateixes probabilitats. \textbf{Sense reemplaçament} no la tornem: en la segona extracció hi ha una bola menys.}

\ex{2} Llancem un dau tres vegades. Quina és la probabilitat de no traure cap 6? I de traure'n almenys un?

\ex{3} Un examen tipus test té 3 preguntes amb 4 opcions cadascuna (només una correcta). Si les contestes totes a l'atzar, quina és la probabilitat d'encertar-les totes? I de no encertar-ne cap? I d'encertar-ne almenys una?

\ex{3} Quantes vegades hem de llançar un dau, com a mínim, perquè la probabilitat de traure almenys un 6 siga més gran que $\frac12$?


% =====================================================
\seccio{6}{Experiències dependents. Probabilitat condicionada}

\ex{1} En una bossa hi ha 5 boles roges i 3 de verdes. Traiem una bola, és roja i no la tornem a la bossa. Quina és la probabilitat que la segona bola siga roja? I que siga verda? Escriu-les amb la notació de la probabilitat condicionada.\recorda{$P(B/A)$ és la probabilitat de $B$ \textbf{sabent que} ha passat $A$; en l'arbre, és la branca de la segona etapa. \ $P(A\cap B)=P(A)\cdot P(B/A)$.}

\ex{1} Traiem dues cartes d'una baralla espanyola de 40 cartes, sense reemplaçament.
\begin{tasks}(1)
  \task Si la primera és un oro, quina és la probabilitat que la segona també ho siga?
  \task Quina és la probabilitat de traure dos oros?
\end{tasks}

\ex{2} En una bossa hi ha 3 boles blanques (B) i 2 de negres (N). En traiem dues \textbf{sense reemplaçament}. Fes el diagrama d'arbre i calcula la probabilitat que:
\begin{tasks}(1)
  \task les dues siguen blanques
  \task les dues siguen negres
  \task n'isca una de cada color. Compara-ho amb l'exercici de les mateixes boles amb reemplaçament
\end{tasks}

\ex{2} En un grup, el 60 \% de l'alumnat aprova Matemàtiques. Dels que aproven Matemàtiques, el 75 \% aprova Física, i dels que no aproven Matemàtiques, només el 25 \%. Fes el diagrama d'arbre i calcula la probabilitat que un alumne triat a l'atzar:
\begin{tasks}(1)
  \task aprove les dues assignatures
  \task aprove Física
  \task no n'aprove cap
\end{tasks}

\ex{3} La urna I té 3 boles roges i 1 de blava; la urna II té 1 bola roja i 3 de blaves. Llancem un dau: si ix 1 o 2, traiem una bola de la urna I, i si no, de la urna II. Quina és la probabilitat de traure una bola roja?

\ex{3} En una bossa hi ha 4 boles roges i 2 de blaves. En traiem tres, una darrere l'altra, sense reemplaçament. Quina és la probabilitat que les tres siguen roges?

\ex{3} En una caixa hi ha 5 bombetes i 2 són defectuoses. En provem dues a l'atzar, sense reemplaçament. Quina és la probabilitat que almenys una siga defectuosa?


% =====================================================
\seccio{7}{Taules de contingència}

\ex{1} En una classe de 25 alumnes hem preguntat qui té mascota. Triem un alumne a l'atzar. Calcula:
\begin{tasks}(1)
  \task $P(\text{té mascota})$
  \task $P(\text{xic i té mascota})$
  \task $P(\text{xica o té mascota})$
  \task si sabem que és una xica, la probabilitat que tinga mascota
\end{tasks}
\par\smallskip\begin{center}{\small\renewcommand{\arraystretch}{1.3}\begin{tabular}{l|cc|c}
\rowcolor{soft} & Mascota & Sense mascota & Total\\\hline
\cellcolor{soft}\textbf{Xiques} & 8 & 4 & \textbf{12}\\
\cellcolor{soft}\textbf{Xics} & 6 & 7 & \textbf{13}\\\hline
\cellcolor{soft}\textbf{Total} & \textbf{14} & \textbf{11} & \textbf{25}\\
\end{tabular}}\end{center}\recorda{En una \textbf{taula de contingència}, els casos possibles són el total (o el total d'una fila, si ja sabem en quina fila està la persona) i els favorables, les caselles que compleixen el que demanen.}

\ex{2} En un institut hi ha 400 alumnes, dels quals 180 són xics. Fan esport federat 120 alumnes, 70 dels quals són xics. Completa la taula i, si triem una persona a l'atzar, calcula:
\begin{tasks}(1)
  \task $P(\text{esport federat})$
  \task $P(\text{xica i esport federat})$
  \task $P(\text{xica o esport federat})$
  \task si sabem que és una xica, la probabilitat que faça esport federat
\end{tasks}
\par\smallskip\begin{center}{\small\renewcommand{\arraystretch}{1.3}\begin{tabular}{l|cc|c}
\rowcolor{soft} & Federat & No federat & Total\\\hline
\cellcolor{soft}\textbf{Xics} & 70 & \phantom{000} & \textbf{180}\\
\cellcolor{soft}\textbf{Xiques} & \phantom{000} & \phantom{000} & \phantom{000}\\\hline
\cellcolor{soft}\textbf{Total} & \textbf{120} & \phantom{000} & \textbf{400}\\
\end{tabular}}\end{center}

\ex{2} En una empresa de 400 treballadors, el 60 \% s'ha vacunat de la grip. Dels vacunats, el 5 \% ha tingut la grip, i dels no vacunats, el 30 \%. Completa la taula i calcula:
\begin{tasks}(1)
  \task la probabilitat que un treballador triat a l'atzar haja tingut la grip
  \task $P(\text{grip}/\text{vacunat})$ i $P(\text{grip}/\text{no vacunat})$
  \task si un treballador ha tingut la grip, la probabilitat que estiga vacunat
\end{tasks}
\par\smallskip\begin{center}{\small\renewcommand{\arraystretch}{1.3}\begin{tabular}{l|cc|c}
\rowcolor{soft} & Grip & Sense grip & Total\\\hline
\cellcolor{soft}\textbf{Vacunats} & \phantom{000} & \phantom{000} & \phantom{000}\\
\cellcolor{soft}\textbf{No vacunats} & \phantom{000} & \phantom{000} & \phantom{000}\\\hline
\cellcolor{soft}\textbf{Total} & \phantom{000} & \phantom{000} & \textbf{400}\\
\end{tabular}}\end{center}\recorda{En una taula de contingència, $P(A/B)=\dfrac{\text{casos de }A\text{ i }B}{\text{casos de }B}=\dfrac{P(A\cap B)}{P(B)}$: només comptem en la fila o la columna de la condició.}

\ex{3} En un grup de 200 persones, el 45 \% són homes i el 30 \% porta ulleres. De les persones amb ulleres, 40 són dones. Fes la taula de contingència i calcula la probabilitat que una persona triada a l'atzar:
\begin{tasks}(1)
  \task siga un home amb ulleres
  \task siga una dona sense ulleres
  \task siga una dona, si sabem que porta ulleres
\end{tasks}

\ex{3} En una classe de 50 alumnes, 30 són de 3r i 20 de 4t. Van a l'institut a peu 12 alumnes de 3r i 8 de 4t. Fes la taula de contingència i compara $P(\text{a peu})$ amb $P(\text{a peu}/\text{3r})$. Són independents els esdeveniments «anar a peu» i «ser de 3r»?

\ex{3} El 70 \% dels clients d'una botiga paga amb targeta. Dels que paguen amb targeta, el 20 \% fa una compra de més de 50 €, i dels que paguen en efectiu, el 10 \%. Fes la taula de contingència per a 1000 clients. Si un client ha fet una compra de més de 50 €, quina és la probabilitat que haja pagat amb targeta?

\vfill
\begin{center}\color{muted}\small\faCheckCircle[regular]\ \ Ànim! Comprova els teus resultats amb les eines interactives dels apunts de la unitat.\end{center}

\end{document}
